This report documents the complete development of a privacy-preserving human
activity recognition system based on mmWave radar point clouds, carried out for
the 23I006 Deep Learning mini project. It records every important decision, the
alternatives considered, the failures met and how each was diagnosed, corrected
and verified. Radar returns are pre-processed with an antenna-return filter,
Doppler-based person centring and spatial gating, which remove about half of all
returns as static clutter. The proposed STPT-Lite applies self-attention within
short temporal segments and across segment tokens and adds an explicit motion
descriptor. It is distilled from a larger teacher, trained with noise-robust augmentation and stored with INT8 weights
(\oursParamsK\,K parameters, \oursMflops\,MFLOPs, \oursLat\,ms per window on one
CPU thread). Under leakage-free protocols with three seeds, it reaches
\rastptlitekdrobustAcc\% on RadHAR and \mistptlitekdrobustAcc\% top-1 on 49 MiliPoint
actions of unseen people, against \mipointnetAcc\% for the best re-implemented baseline. An automated
audit shows that random window splits place most test frames inside training
windows, which explains the large spread of published results. The report also
covers the robustness study, ablations, cross-dataset transfer, explainability,
the working demonstration, testing and the remaining limitations.
